\documentclass[a4paper,14pt]{extarticle}

% ------------------------------------------------------------
%  Пакеты и базовые настройки
% ------------------------------------------------------------
\usepackage[T2A]{fontenc}
\usepackage[utf8]{inputenc}
\usepackage[russian]{babel}
\usepackage{tempora}

\usepackage{geometry}
\geometry{
  left=30mm,
  right=20mm,
  top=20mm,
  bottom=20mm,
  includeheadfoot
}

\usepackage{setspace}
\onehalfspacing

\usepackage{graphicx}
\usepackage{float}
\usepackage{caption}
\captionsetup[table]{justification=raggedright,singlelinecheck=false,labelsep=endash}

\usepackage{array}
\usepackage{tabularx}
\usepackage{longtable}
\usepackage{multirow}
\usepackage{makecell}

\usepackage{enumitem}
\usepackage{indentfirst}
\usepackage{fancyhdr}
\usepackage{lastpage}
\usepackage{microtype}
\usepackage{ragged2e}
\usepackage{seqsplit}
\usepackage{hyperref}

\hypersetup{
  colorlinks=false,
  pdfborder={0 0 0}
}

% ------------------------------------------------------------
%  Абзацы, списки, колонтитулы
% ------------------------------------------------------------
\setlength{\parindent}{1.25cm}
\setlength{\parskip}{0pt}
\setlist{nosep,leftmargin=*}

\pagestyle{fancy}
\fancyhf{}
\rfoot{Страница \thepage{} из \pageref{LastPage}}
\renewcommand{\headrulewidth}{0pt}

% ------------------------------------------------------------
%  Макросы и таблицы
% ------------------------------------------------------------
\newcommand{\TBD}{\textbf{[TBD]}}
\newcommand{\OrgName}{Кафедра компьютерных систем и технологий НИЯУ МИФИ (№12)}
\newcommand{\SystemFullName}{Автоматизированная мидл-офис система для ведения позиций по счетам ностро и операционной отчетности}
\newcommand{\SystemShortName}{FX-АСУБАНК}
\newcommand{\DocCipher}{ШИФР-ПЗТП}
\newcommand{\DocVersion}{1.0}
\newcommand{\DocYear}{2026}

\setlength{\tabcolsep}{4pt}
\renewcommand{\arraystretch}{1.12}
\providecommand{\textslash}{/\penalty\hyphenpenalty\relax}
\let\oldtexttt\texttt
\renewcommand{\texttt}[1]{\oldtexttt{\seqsplit{#1}}}

\newcolumntype{L}[1]{>{\sloppy\RaggedRight\let\newline\\\arraybackslash\hspace{0pt}}p{#1}}
\newcolumntype{C}[1]{>{\sloppy\Centering\let\newline\\\arraybackslash\hspace{0pt}}p{#1}}
\newcolumntype{R}[1]{>{\sloppy\RaggedLeft\let\newline\\\arraybackslash\hspace{0pt}}p{#1}}

\begin{document}
\sloppy

% ============================================================
%  Титульный лист
% ============================================================
\thispagestyle{empty}
\begin{center}
\textbf{\MakeUppercase{\OrgName}}
\end{center}

\vspace{28mm}

\begin{center}
\textbf{\MakeUppercase{\SystemFullName}}\\[6mm]
\textbf{\MakeUppercase{ПОЯСНИТЕЛЬНАЯ ЗАПИСКА\\К ТЕХНИЧЕСКОМУ ПРОЕКТУ}}\\[4mm]
\DocCipher\\[4mm]
Листов \pageref{LastPage}\\[4mm]
Версия \DocVersion
\end{center}

\vfill

\begin{center}
Москва, \DocYear
\end{center}

\newpage

% ============================================================
%  Аннотация
% ============================================================
\begin{center}
\textbf{\MakeUppercase{АННОТАЦИЯ}}
\end{center}

В настоящей пояснительной записке к техническому проекту представлены автоматизируемый процесс проведения валютных сделок, основные технические решения по созданию автоматизированной системы \SystemShortName, решения по составу подсистем, информационному обмену, программному и техническому обеспечению, а также мероприятия по подготовке объекта автоматизации к вводу системы в действие.

Документ разработан на основании технического задания на систему \SystemShortName\ и с учетом структуры пояснительной записки к техническому проекту, предусмотренной требованиями комплекса стандартов на автоматизированные системы.

\newpage

% ============================================================
%  Содержание
% ============================================================
\tableofcontents
\newpage

% ============================================================
%  Термины и определения
% ============================================================
\begin{center}
\textbf{\MakeUppercase{ТЕРМИНЫ И ОПРЕДЕЛЕНИЯ}}
\end{center}

\begin{description}[leftmargin=1.5cm,itemsep=1mm]
  \item[Архитектура] Совокупность подсистем, компонентов, интерфейсов, принципов взаимодействия и развертывания автоматизированной системы.
  \item[Валютная сделка] Сделка купли или продажи одной валюты за другую валюту.
  \item[Дата валютирования] Дата фактического исполнения расчетов по валютной сделке.
  \item[Мидл-офис] Подразделение банка, выполняющее контроль, согласование и подтверждение сделок.
  \item[Ностро-счет] Счет банка, открытый в другой кредитной организации и используемый для расчетов в соответствующей валюте.
  \item[Позиция] Состояние требований и обязательств банка по валюте, счету или контрагенту.
  \item[Подсистема] Часть автоматизированной системы, выделенная по функциональному или технологическому признаку.
  \item[Технический проект] Проектная стадия, на которой фиксируются основные технические решения по созданию автоматизированной системы.
\end{description}

\newpage

\begin{center}
\textbf{\MakeUppercase{ОБОЗНАЧЕНИЯ И СОКРАЩЕНИЯ}}
\end{center}

\begin{table}[H]
\centering
\small
\caption{Обозначения и сокращения}
\begin{tabularx}{\linewidth}{|L{35mm}|X|}
\hline
\textbf{Сокращение} & \textbf{Расшифровка} \\
\hline
АС & Автоматизированная система \\
\hline
БД & База данных \\
\hline
ИБ & Информационная безопасность \\
\hline
НСИ & Нормативно-справочная информация \\
\hline
ПЗ & Пояснительная записка \\
\hline
ПО & Программное обеспечение \\
\hline
ПМИ & Программа и методика испытаний \\
\hline
СУБД & Система управления базами данных \\
\hline
ТП & Технический проект \\
\hline
ТЗ & Техническое задание \\
\hline
API & Application Programming Interface \\
\hline
FX & Foreign Exchange, валютные операции \\
\hline
P\&L & Profit and Loss, финансовый результат \\
\hline
REST & Representational State Transfer \\
\hline
SLA & Service Level Agreement \\
\hline
TLS & Transport Layer Security \\
\hline
UAT & User Acceptance Testing, приемочные испытания пользователем \\
\hline
\end{tabularx}
\end{table}

\newpage

% ============================================================
%  1 Общие положения
% ============================================================
\section{Общие положения}

\subsection{Наименование проектируемой Системы и ее условное обозначение}
Полное наименование проектируемой системы: \SystemFullName.

Условное обозначение проектируемой системы: \SystemShortName.

Далее по тексту используется также обозначение «Система».

\subsection{Перечень документов, на основании которых проектируется Система}
Техническое проектирование Системы выполняется на основании следующих документов:
\begin{itemize}[itemsep=1mm]
  \item техническое задание на создание \SystemShortName;
  \item РУП 09.05.01 по дисциплине «Применение и эксплуатация автоматизированных систем специального назначения НИЯУ МИФИ»;
  \item ГОСТ 34.601-90 «Автоматизированные системы. Стадии создания»;
  \item ГОСТ Р 59795-2021 «Информационные технологии. Комплекс стандартов на автоматизированные системы. Автоматизированные системы. Требования к содержанию документов»;
  \item ГОСТ 34.201-89 «Виды, комплектность и обозначение документов при создании автоматизированных систем».
\end{itemize}

\subsection{Перечень организаций, участвующих в разработке Системы}
Предприятие-заказчик: \OrgName.

Предприятие-разработчик: команда №1 «Сделки FX» в рамках учебного проекта.

Смежные участники проекта:
\begin{itemize}[itemsep=1mm]
  \item проект системы ведения позиций по счетам банка;
  \item проект системы формирования отчетов и ведения НСИ.
\end{itemize}

\subsection{Плановые сроки начала работы по созданию Системы}
Плановый срок начала работ по созданию Системы: 24.02.2026.

Плановый срок окончания работ по созданию Системы: не позднее 25.05.2026.

\subsection{Цели создания Системы}
Цели создания Системы и критерии оценки достижения целей приведены в таблице \ref{tab:goals}.

\begin{table}[H]
\centering
\small
\caption{Цели Системы и критерии оценки достижения целей}
\label{tab:goals}
\begin{tabularx}{\linewidth}{|L{46mm}|L{40mm}|X|}
\hline
\textbf{Цель} & \textbf{Показатель} & \textbf{Критерии оценки достижения целей} \\
\hline
Автоматизация жизненного цикла FX-сделки от фиксации намерения до расчетов &
Уровень автоматизации ключевых этапов &
Реализованы создание, проверка, согласование, расчет платежей, подбор ностро-счетов и передача данных для обновления позиций. \\
\hline
Снижение ошибок ручного ввода и согласования &
Количество ошибок при внесении данных &
Не более 1 ошибки на 10000 операций за счет использования НСИ, форматно-логических проверок и ролевого контроля. \\
\hline
Ускорение контроля параметров сделки &
Время выполнения проверки &
Автоматическая проверка параметров сделки выполняется не более чем за 1 секунду при штатной нагрузке. \\
\hline
Повышение прослеживаемости операций &
Полнота аудита &
Все изменения статусов и ключевых полей сделки фиксируются с указанием пользователя, времени, старого и нового значения. \\
\hline
\end{tabularx}
\end{table}

\subsection{Назначение Системы}
Система предназначена для автоматизации банковского процесса проведения валютных сделок типов TOD, TOM и Spot в коммерческом банке: от регистрации параметров сделки трейдером до согласования мидл-офисом, формирования расчетных платежей, обновления позиций по счетам ностро и подготовки операционной отчетности.

\subsection{Подтверждение соответствия проектных решений действующим нормам и правилам техники безопасности, пожаро- и взрывобезопасности}
Проектные решения предусматривают эксплуатацию Системы на типовой серверной и пользовательской инфраструктуре Заказчика. Требования техники безопасности, электробезопасности, пожарной безопасности и санитарные требования обеспечиваются инфраструктурными службами объекта автоматизации и применяемыми регламентами эксплуатации оборудования.

Технические средства, используемые для размещения Системы, должны эксплуатироваться в помещениях, оснащенных защитным заземлением, средствами пожарной сигнализации, источниками бесперебойного питания для критичных серверных компонентов и средствами контроля параметров окружающей среды.

\subsection{Сведения об использованных при проектировании нормативно-технических документах}
При разработке технического проекта используются следующие нормативно-технические документы:
\begin{itemize}[itemsep=1mm]
  \item ГОСТ 34.601-90 «Автоматизированные системы. Стадии создания»;
  \item ГОСТ 34.201-89 «Информационная технология. Комплекс стандартов на автоматизированные системы. Виды, комплектность и обозначение документов при создании автоматизированных систем»;
  \item ГОСТ Р 59795-2021 «Информационные технологии. Комплекс стандартов на автоматизированные системы. Автоматизированные системы. Требования к содержанию документов»;
  \item ГОСТ 19.201-78 «Единая система программной документации. Техническое задание. Требования к содержанию и оформлению»;
  \item ГОСТ 2.105-2019 «Единая система конструкторской документации. Общие требования к текстовым документам».
\end{itemize}

\subsection{Сведения о НИР, передовом опыте, изобретениях, использованных при разработке проекта}
При разработке проекта используются типовые подходы к построению банковских мидл-офисных систем: разделение функций фронт-офиса и мидл-офиса, централизованное ведение НСИ, событийная передача статусов сделок, журналирование действий пользователей, ролевое разграничение доступа и интеграция с системами позиций и отчетности.

Сведения об использовании изобретений, патентов или результатов научно-исследовательских работ отсутствуют.

\subsection{Очередность создания Системы и объем каждой очереди}
Создание Системы выполняется одной очередью в составе минимально необходимого функционального контура. В рамках очереди предусматриваются:
\begin{itemize}[itemsep=1mm]
  \item разработка проектной документации;
  \item реализация подсистемы ввода и ведения FX-сделок;
  \item реализация расчетной логики и проверки параметров сделки;
  \item реализация аудита, журналирования и ролевой модели;
  \item реализация интеграционных интерфейсов со смежными системами;
  \item подготовка демонстрационных и приемочных испытаний.
\end{itemize}

% ============================================================
%  2 Описание процесса деятельности
% ============================================================
\section{Описание процесса деятельности}

Автоматизируемый процесс относится к деятельности фронт-офиса, мидл-офиса и операционного блока банка при проведении валютных операций. До внедрения Системы отдельные действия могут выполняться вручную либо с использованием разрозненных инструментов, что повышает риск ошибок при вводе параметров сделки, согласовании и формировании расчетной информации.

Проектируемый процесс деятельности включает следующие этапы:
\begin{enumerate}[label=\arabic*.,itemsep=1mm]
  \item трейдер регистрирует параметры FX-сделки: валютную пару, сумму, курс, тип сделки, контрагента и дату сделки;
  \item Система выполняет форматно-логический контроль и проверку НСИ;
  \item Система рассчитывает дату валютирования и платежи по сделке;
  \item сделка передается на согласование мидл-офису;
  \item позиционер принимает решение об одобрении, возврате на доработку или отклонении сделки;
  \item при одобрении данные сделки передаются в систему ведения позиций по счетам банка;
  \item сведения о сделках используются для формирования отчетности и контроля операционного дня.
\end{enumerate}

Организация работ в условиях функционирования Системы должна обеспечивать разделение полномочий между пользователями. Трейдер отвечает за ввод и корректность исходных коммерческих параметров сделки. Позиционер отвечает за контроль, подтверждение и обработку исключительных ситуаций. Администратор отвечает за учетные записи, роли, настройки и контроль справочников.

\begin{table}[H]
\centering
\small
\caption{Роли участников автоматизируемого процесса}
\label{tab:roles}
\begin{tabularx}{\linewidth}{|L{35mm}|L{45mm}|X|}
\hline
\textbf{Роль} & \textbf{Зона ответственности} & \textbf{Основные действия в Системе} \\
\hline
Трейдер & Проведение валютных сделок & Создание черновика сделки, редактирование параметров, передача на согласование, просмотр статуса. \\
\hline
Позиционер & Контроль и подтверждение сделок & Проверка параметров, одобрение, отклонение, возврат на доработку, контроль расчетных данных. \\
\hline
Аналитик/руководитель & Контроль операций и отчетность & Просмотр реестра сделок, формирование отчетов, анализ статусов и операционных показателей. \\
\hline
Администратор & Эксплуатационное сопровождение & Управление пользователями, ролями, настройками, контроль журналов и работоспособности. \\
\hline
\end{tabularx}
\end{table}

% ============================================================
%  3 Основные технические решения
% ============================================================
\section{Основные технические решения}

\subsection{Ограничения на технические решения}
При проектировании Системы принимаются следующие ограничения:
\begin{itemize}[itemsep=1mm]
  \item Система создается как клиент-серверное веб-приложение;
  \item хранение оперативных данных выполняется в централизованной реляционной БД;
  \item обмен со смежными системами выполняется через согласованные API или асинхронные сообщения;
  \item проектные решения должны обеспечивать трассируемость требований ТЗ до функций и приемочных критериев;
  \item интерфейс пользователя должен быть русскоязычным;
  \item технические решения должны позволять развертывание в локальном ЦОД или защищенном облачном окружении Заказчика.
\end{itemize}

\subsection{Решения по структуре Системы, подсистем и информационному обмену между компонентами}
В состав Системы включаются подсистемы, приведенные в таблице \ref{tab:subsystems}.

\begin{table}[H]
\centering
\small
\caption{Состав подсистем Системы}
\label{tab:subsystems}
\begin{tabularx}{\linewidth}{|L{45mm}|X|}
\hline
\textbf{Подсистема} & \textbf{Назначение} \\
\hline
Подсистема ввода и ведения FX-сделок & Создание, редактирование, отмена, поиск и просмотр сделок, управление статусами сделки в рамках разрешенных переходов. \\
\hline
Подсистема расчетов & Расчет даты валютирования, платежей, направления списания/зачисления, подбор расчетных ностро-счетов. \\
\hline
Подсистема согласования & Передача сделки на контроль мидл-офису, фиксация решения, обработка возврата и отклонения. \\
\hline
Подсистема аудита и журналирования & Регистрация действий пользователей, изменений статусов, обмена со смежными системами и ошибок обработки. \\
\hline
Подсистема администрирования & Управление пользователями, ролями, настройками, справочниками и параметрами доступа. \\
\hline
Подсистема взаимодействия & Обмен данными с системами позиций, отчетности и НСИ. \\
\hline
\end{tabularx}
\end{table}

\subsubsection{Схема компонент/модулей Системы}
Логическая схема компонентов Системы представлена в таблице \ref{tab:components}. Графическое представление схемы компонентов подлежит включению в приложение при наличии утвержденной UML-диаграммы.

\begin{table}[H]
\centering
\small
\caption{Логические компоненты Системы}
\label{tab:components}
\begin{tabularx}{\linewidth}{|L{40mm}|L{35mm}|X|}
\hline
\textbf{Компонент} & \textbf{Тип} & \textbf{Основные функции} \\
\hline
Веб-клиент & Клиентский компонент & Экранные формы для трейдера, позиционера, аналитика и администратора. \\
\hline
Сервис сделок & Серверный компонент & Управление жизненным циклом сделки, проверка переходов статусов, хранение карточки сделки. \\
\hline
Расчетный сервис & Серверный компонент & Расчет даты валютирования, платежей, направления движения денежных средств и ностро-счетов. \\
\hline
Сервис НСИ & Серверный компонент & Получение и кэширование справочников валют, контрагентов, счетов и календарей. \\
\hline
Интеграционный модуль & Серверный компонент & Передача и прием сообщений от смежных систем, повторные попытки, контроль ошибок обмена. \\
\hline
Журнал аудита & Инфраструктурный компонент & Хранение событий действий пользователей, изменений данных и интеграционных взаимодействий. \\
\hline
База данных & Инфраструктурный компонент & Централизованное хранение оперативных, справочных и аудиторских данных. \\
\hline
\end{tabularx}
\end{table}

\subsection{Решения по взаимосвязям Системы со смежными системами и обеспечению совместимости}
Система взаимодействует со смежными системами по направлениям, приведенным в таблице \ref{tab:integrations}.

{\centering
\small
\setlength{\tabcolsep}{3pt}
\begin{longtable}{|L{31mm}|L{28mm}|L{31mm}|L{26mm}|L{23mm}|}
\caption{Интеграционные взаимодействия}\label{tab:integrations}\\
\hline
\textbf{Смежная система} & \textbf{Направление} & \textbf{Состав данных} & \textbf{Режим обмена} & \textbf{Требования} \\
\hline
\endfirsthead
\hline
\textbf{Смежная система} & \textbf{Направление} & \textbf{Состав данных} & \textbf{Режим обмена} & \textbf{Требования} \\
\hline
\endhead
\hline
\multicolumn{5}{|r|}{\textit{Продолжение на след. странице}} \\
\hline
\endfoot
\hline
\endlastfoot

Система ведения позиций & \SystemShortName{} $\rightarrow$ смежная система & ID сделки, валютная пара, суммы платежей, контрагент, дата валютирования, статус & По событию одобрения сделки & Гарантированная доставка, подтверждение приема, журналирование \\
\hline
Система ведения позиций & Смежная система $\rightarrow$ \SystemShortName{} & Решение, комментарий, дата подтверждения, идентификатор позиции & По событию решения & Идемпотентность обработки, контроль повторов \\
\hline
Система отчетности и НСИ & Смежная система $\rightarrow$ \SystemShortName{} & Справочники валют, контрагентов, счетов ностро, производственный календарь & Регламентно или по событию & Контроль версии справочников, валидация состава данных \\
\hline
Система отчетности и НСИ & \SystemShortName{} $\rightarrow$ смежная система & Данные по сделкам и статусам для отчетов & По запросу или в конце операционного дня & Согласованный формат выгрузки, контроль полноты \\
\hline
\end{longtable}
\par}

Для синхронных интерфейсов допускается применение REST API с JSON. Для асинхронных взаимодействий допускается применение очередей сообщений. Все интеграционные запросы и ответы должны протоколироваться с фиксацией времени, идентификатора корреляции, результата обработки и текста ошибки при ее наличии.

\subsection{Решения по режимам функционирования и диагностированию работы Системы}
Система должна поддерживать следующие режимы функционирования:
\begin{itemize}[itemsep=1mm]
  \item полнофункциональный режим, в котором доступны операции создания, согласования, расчета, интеграционного обмена и отчетности;
  \item режим ограниченной функциональности, применяемый при регламентных работах, обновлениях и обслуживании;
  \item режим восстановления после сбоя, в котором выполняется проверка целостности данных, повтор интеграционных сообщений и контроль незавершенных операций.
\end{itemize}

Диагностирование работы Системы обеспечивается следующими средствами:
\begin{itemize}[itemsep=1mm]
  \item журнал ошибок прикладных компонентов;
  \item журнал аудита действий пользователей;
  \item контроль доступности БД и внешних интерфейсов;
  \item контроль очередей интеграционных сообщений;
  \item регистрация метрик производительности и времени отклика.
\end{itemize}

\subsection{Решения по численности, квалификации и функциям персонала АС}
Для эксплуатации Системы предусматриваются категории персонала, приведенные в таблице \ref{tab:staff}.

\begin{table}[H]
\centering
\small
\caption{Персонал Системы}
\label{tab:staff}
\begin{tabularx}{\linewidth}{|L{35mm}|L{45mm}|X|}
\hline
\textbf{Категория} & \textbf{Квалификация} & \textbf{Функции} \\
\hline
Пользователь-трейдер & Знание процесса FX-сделок и правил оформления операций & Ввод и сопровождение сделок, исправление замечаний мидл-офиса. \\
\hline
Пользователь-позиционер & Знание правил контроля сделок, расчетов и позиций & Проверка сделок, согласование, возврат на доработку, контроль расчетных данных. \\
\hline
Администратор & Навыки администрирования пользователей, ролей и прикладных настроек & Управление доступом, справочниками, мониторингом и журналами. \\
\hline
Технический специалист & Навыки сопровождения серверного ПО, СУБД и интеграций & Развертывание, резервное копирование, восстановление, контроль работоспособности. \\
\hline
\end{tabularx}
\end{table}

\subsection{Сведения об обеспечении заданных в ТЗ потребительских характеристик Системы}
Ключевые потребительские характеристики Системы обеспечиваются решениями, приведенными в таблице \ref{tab:quality}.

\begin{table}[H]
\centering
\small
\caption{Обеспечение потребительских характеристик}
\label{tab:quality}
\begin{tabularx}{\linewidth}{|L{35mm}|L{45mm}|X|}
\hline
\textbf{Характеристика} & \textbf{Требуемый результат} & \textbf{Проектное решение} \\
\hline
Производительность & Проверка параметров сделки не более 1 секунды & Локальные проверки формата, индексы БД, кэширование НСИ и календарей, минимизация синхронных внешних вызовов. \\
\hline
Надежность & Сохранность данных сделки при отказах & Транзакционная запись операции и платежей, повтор интеграционных сообщений, резервное копирование БД. \\
\hline
Безопасность & Разграничение доступа и аудит & Ролевая модель, аутентификация, авторизация операций, неизменяемый журнал аудита. \\
\hline
Качество данных & Снижение ошибок ручного ввода & Использование справочников валют, контрагентов, счетов; форматно-логические проверки; контроль активного статуса контрагента. \\
\hline
Сопровождаемость & Возможность диагностики и развития & Разделение на подсистемы, журналирование, мониторинг, документированные интерфейсы. \\
\hline
\end{tabularx}
\end{table}

\subsection{Состав функций, реализуемых Системой}
Состав функций Системы приведен в таблице \ref{tab:functions}. Полная детализация функциональных требований приведена в техническом задании.

{\centering
\small
\setlength{\tabcolsep}{3pt}
\begin{longtable}{|C{10mm}|L{45mm}|L{32mm}|L{48mm}|}
\caption{Функции Системы}\label{tab:functions}\\
\hline
\textbf{№} & \textbf{Функция} & \textbf{Пользователь} & \textbf{Входные/выходные данные} \\
\hline
\endfirsthead
\hline
\textbf{№} & \textbf{Функция} & \textbf{Пользователь} & \textbf{Входные/выходные данные} \\
\hline
\endhead
\hline
\multicolumn{4}{|r|}{\textit{Продолжение на след. странице}} \\
\hline
\endfoot
\hline
\endlastfoot

1 & Создание FX-сделки & Трейдер & Вход: параметры сделки. Выход: карточка сделки в статусе «Черновик». \\
\hline
2 & Проверка параметров сделки & Система & Вход: карточка сделки, НСИ, календарь. Выход: результат проверки, список ошибок. \\
\hline
3 & Расчет платежей и даты валютирования & Система & Вход: тип сделки, дата сделки, валютная пара, сумма, курс. Выход: платежи и дата валютирования. \\
\hline
4 & Передача сделки на согласование & Трейдер & Вход: проверенная сделка. Выход: задача мидл-офису, статус «На согласовании». \\
\hline
5 & Одобрение или возврат сделки & Позиционер & Вход: карточка сделки, комментарий. Выход: решение и новый статус сделки. \\
\hline
6 & Передача данных в систему позиций & Система & Вход: одобренная сделка. Выход: интеграционное сообщение и результат обработки. \\
\hline
7 & Формирование отчетов & Аналитик/руководитель & Вход: период, фильтры. Выход: отчет по сделкам и статусам. \\
\hline
8 & Администрирование доступа & Администратор & Вход: пользователи, роли, настройки. Выход: права доступа и журнал изменений. \\
\hline
\end{longtable}
\par}

\subsection{Решения по комплексу технических средств и его размещению на объекте}
Система размещается на серверной инфраструктуре Заказчика либо в защищенном облачном окружении. Минимальный состав технических средств приведен в таблице \ref{tab:hardware}.

\begin{table}[H]
\centering
\small
\caption{Минимальный состав технических средств}
\label{tab:hardware}
\begin{tabularx}{\linewidth}{|L{40mm}|L{45mm}|X|}
\hline
\textbf{Средство} & \textbf{Минимальные характеристики} & \textbf{Назначение} \\
\hline
Сервер приложений & 4 vCore, 8 ГБ RAM, SSD 100 ГБ & Размещение серверных компонентов и API. \\
\hline
Сервер БД & 4 vCore, 16 ГБ RAM, SSD 200 ГБ, RAID 10 или аналог & Хранение оперативных, справочных и аудиторских данных. \\
\hline
Сервер кэша/очередей & 2 vCore, 4 ГБ RAM, SSD 50 ГБ & Кэширование НСИ, хранение сессий, обработка асинхронных сообщений. \\
\hline
Рабочие станции пользователей & Современный браузер Chrome/Firefox/Edge & Работа пользователей с веб-интерфейсом. \\
\hline
Сетевое оборудование & LAN не менее 1 Гбит/с, защищенные каналы связи & Доступ пользователей и интеграционный обмен. \\
\hline
\end{tabularx}
\end{table}

\subsection{Решения по составу информации, объему, способам организации, входным и выходным документам}
Информационное обеспечение Системы строится на централизованной реляционной БД. Данные должны быть организованы в соответствии с логической моделью, включающей сущности сделки, платежа, контрагента, валюты, ностро-счета, календаря, пользователя, роли и события аудита.

Основные информационные объекты приведены в таблице \ref{tab:data}.

\begin{table}[H]
\centering
\small
\caption{Основные информационные объекты}
\label{tab:data}
\begin{tabularx}{\linewidth}{|L{35mm}|L{50mm}|X|}
\hline
\textbf{Объект} & \textbf{Ключевые атрибуты} & \textbf{Назначение} \\
\hline
Сделка & ID, тип, валютная пара, сумма, курс, контрагент, статус, дата сделки & Учет жизненного цикла FX-сделки. \\
\hline
Платеж & ID, сделка, валюта, сумма, направление, счет, дата валютирования & Описание расчетных движений по сделке. \\
\hline
Контрагент & ID, наименование, статус активности, реквизиты & Проверка допустимости сделки и формирование отчетности. \\
\hline
Валюта & Код, наименование, точность, статус & Расчет и контроль валютных параметров. \\
\hline
Ностро-счет & ID, валюта, банк-корреспондент, статус & Подбор расчетных счетов для платежей. \\
\hline
Событие аудита & ID, пользователь, действие, объект, время, старое/новое значение & Прослеживаемость действий и изменений. \\
\hline
\end{tabularx}
\end{table}

Входная информация включает параметры сделки, данные НСИ, производственный календарь и решения мидл-офиса. Выходная информация включает карточки сделок, расчетные платежи, интеграционные сообщения в систему позиций, отчеты по сделкам и журналы аудита.

\subsection{Решения по составу программных средств, языкам программирования, алгоритмам и методам реализации}
Программное обеспечение Системы включает системное и прикладное ПО. Минимальный состав приведен в таблице \ref{tab:software}.

\begin{table}[H]
\centering
\small
\caption{Состав программных средств}
\label{tab:software}
\begin{tabularx}{\linewidth}{|L{45mm}|X|}
\hline
\textbf{Категория} & \textbf{Программные средства} \\
\hline
Операционная система & Linux Ubuntu 20.04+/RHEL 8+ или Windows Server 2019+. \\
\hline
СУБД & PostgreSQL версии 14 и выше. \\
\hline
Кэширование и очереди & Redis 6+ и/или брокер сообщений \texttt{RabbitMQ\textslash Kafka}. \\
\hline
Серверная часть & Python 3.10+ или Node.js 18+ в зависимости от утвержденного стека реализации. \\
\hline
Клиентская часть & Веб-приложение, работающее в Chrome 90+, Firefox 88+, Edge 90+. \\
\hline
Средства поставки & Docker 20+ и Docker Compose при контейнерном развертывании. \\
\hline
\end{tabularx}
\end{table}

Алгоритмы расчета даты валютирования должны учитывать тип сделки: TOD, TOM или Spot, а также производственный календарь. Алгоритмы проверки параметров должны обеспечивать контроль положительности суммы, различия валют валютной пары, активного статуса контрагента и наличия расчетных счетов.

\subsection{Решения по обеспечению информационной безопасности}
Основные угрозы информационной безопасности Системы:
\begin{itemize}[itemsep=1mm]
  \item несанкционированный доступ к карточкам сделок;
  \item изменение параметров сделки без полномочий;
  \item подмена или потеря интеграционных сообщений;
  \item утечка справочной или операционной информации;
  \item нарушение целостности данных при отказах.
\end{itemize}

Для снижения указанных угроз предусматриваются следующие решения:
\begin{itemize}[itemsep=1mm]
  \item идентификация и аутентификация пользователей;
  \item ролевая модель доступа с разделением функций трейдера, позиционера, аналитика и администратора;
  \item применение TLS для защищенного обмена данными;
  \item аудит входа, выхода, изменений сделок, решений и интеграционных операций;
  \item контроль целостности данных средствами СУБД;
  \item резервное копирование и регламент восстановления данных.
\end{itemize}

% ============================================================
%  4 Мероприятия по подготовке объекта автоматизации
% ============================================================
\section{Мероприятия по подготовке объекта автоматизации к вводу Системы в действие}

\subsection{Мероприятия по приведению информации к виду, пригодному для обработки на ЭВМ}
До ввода Системы в действие необходимо подготовить и проверить справочники валют, контрагентов, ностро-счетов, пользователей, ролей и производственный календарь. Данные НСИ должны быть очищены от дублей, неактуальных записей и противоречивых статусов.

\begin{table}[H]
\centering
\small
\caption{Мероприятия подготовки данных}
\label{tab:data-prep}
\begin{tabularx}{\linewidth}{|L{45mm}|L{45mm}|X|}
\hline
\textbf{Мероприятие} & \textbf{Ответственный} & \textbf{Результат} \\
\hline
Подготовка справочника валют & Владелец НСИ & Загружен актуальный перечень валют с кодами и точностью. \\
\hline
Подготовка справочника контрагентов & Владелец НСИ/мидл-офис & Загружены активные контрагенты, указаны статусы и идентификаторы. \\
\hline
Подготовка справочника счетов ностро & Операционный блок & Загружены счета, валюты, банки-корреспонденты и статусы доступности. \\
\hline
Подготовка пользователей и ролей & Администратор & Созданы учетные записи и назначены роли. \\
\hline
Проверка тестового набора сделок & Разработчик/Заказчик & Подтверждена корректность расчетов и статусов на контрольных примерах. \\
\hline
\end{tabularx}
\end{table}

\subsection{Мероприятия по обучению и проверке квалификации персонала}
Перед вводом Системы в действие необходимо провести обучение пользователей и администраторов. Программа обучения должна включать:
\begin{itemize}[itemsep=1mm]
  \item общую концепцию работы Системы;
  \item порядок создания и редактирования FX-сделки;
  \item порядок согласования, возврата и отклонения сделки;
  \item работу с реестром сделок и отчетами;
  \item действия при ошибках и нештатных ситуациях;
  \item администрирование пользователей, ролей и справочников.
\end{itemize}

Проверка квалификации выполняется путем демонстрационного выполнения типовых сценариев в тестовом контуре.

\subsection{Мероприятия по созданию необходимых подразделений и рабочих мест}
Создание новых подразделений для эксплуатации Системы не требуется. Рабочие места пользователей должны быть оснащены персональными компьютерами с доступом к корпоративной сети и поддерживаемым веб-браузером.

Для администраторов и технических специалистов должен быть обеспечен доступ к средствам мониторинга, журналам, инструментам резервного копирования и регламентным процедурам сопровождения.

\subsection{Мероприятия по изменению объекта автоматизации}
Ввод Системы требует актуализации регламентов работы фронт-офиса и мидл-офиса. В регламентах должны быть зафиксированы:
\begin{itemize}[itemsep=1mm]
  \item порядок регистрации FX-сделки;
  \item правила перевода сделки между статусами;
  \item порядок контроля и согласования мидл-офисом;
  \item ответственность за ведение НСИ;
  \item порядок обработки ошибок интеграционного обмена;
  \item порядок формирования и использования отчетности.
\end{itemize}

По завершении подготовительных мероприятий проводится демонстрация готовности Системы и оформляется результат приемки в установленном порядке.

% ============================================================
%  Список источников
% ============================================================
\newpage
\begin{center}
\textbf{\MakeUppercase{СПИСОК ИСПОЛЬЗОВАННЫХ ИСТОЧНИКОВ}}
\end{center}

\begin{enumerate}[label=\arabic*.,itemsep=1mm]
  \item ГОСТ 34.601-90. Информационная технология. Комплекс стандартов на автоматизированные системы. Автоматизированные системы. Стадии создания.
  \item ГОСТ 34.201-89. Информационная технология. Комплекс стандартов на автоматизированные системы. Виды, комплектность и обозначение документов при создании автоматизированных систем.
  \item ГОСТ Р 59795-2021. Информационные технологии. Комплекс стандартов на автоматизированные системы. Автоматизированные системы. Требования к содержанию документов.
  \item ГОСТ 2.105-2019. Единая система конструкторской документации. Общие требования к текстовым документам.
  \item Техническое задание на создание автоматизированной мидл-офис системы для ведения позиций по счетам ностро и операционной отчетности.
\end{enumerate}

% ============================================================
%  Приложения
% ============================================================
\newpage
\begin{center}
\textbf{Приложение А. Схема компонентов Системы}
\end{center}

Схема компонентов Системы представляется следующей логической структурой:
\begin{itemize}[itemsep=1mm]
  \item пользовательские рабочие места взаимодействуют с веб-клиентом;
  \item веб-клиент обращается к серверному API;
  \item серверный API маршрутизирует запросы в сервис сделок, расчетный сервис, сервис НСИ, сервис аудита и интеграционный модуль;
  \item серверные компоненты используют единую БД и кэш справочной информации;
  \item интеграционный модуль взаимодействует со смежными системами позиций, отчетности и НСИ.
\end{itemize}

\begin{table}[H]
\centering
\small
\caption{Связи компонентов Системы}
\begin{tabularx}{\linewidth}{|L{40mm}|L{40mm}|X|}
\hline
\textbf{Источник} & \textbf{Получатель} & \textbf{Передаваемая информация} \\
\hline
Веб-клиент & Серверный API & Пользовательские запросы, параметры сделок, команды согласования. \\
\hline
Серверный API & Сервис сделок & Команды создания, изменения, поиска и изменения статуса сделки. \\
\hline
Сервис сделок & Расчетный сервис & Параметры сделки для расчета даты валютирования и платежей. \\
\hline
Сервис сделок & Сервис аудита & События изменения статуса и параметров сделки. \\
\hline
Интеграционный модуль & Смежные системы & Сообщения по сделкам, справочникам, отчетным данным и решениям. \\
\hline
\end{tabularx}
\end{table}

\newpage
\begin{center}
\textbf{Приложение Б. Схема пользовательского интерфейса}
\end{center}

Пользовательский интерфейс Системы должен включать следующие основные экранные формы:
\begin{itemize}[itemsep=1mm]
  \item форма авторизации;
  \item реестр сделок;
  \item карточка FX-сделки;
  \item форма создания/редактирования сделки;
  \item рабочее место позиционера для согласования;
  \item журнал аудита;
  \item раздел отчетов;
  \item раздел администрирования пользователей, ролей и справочников.
\end{itemize}

\begin{table}[H]
\centering
\small
\caption{Экранные формы и доступные роли}
\begin{tabularx}{\linewidth}{|L{45mm}|L{45mm}|X|}
\hline
\textbf{Экранная форма} & \textbf{Основные действия} & \textbf{Роли} \\
\hline
Реестр сделок & Поиск, фильтрация, просмотр карточки & Трейдер, позиционер, аналитик, администратор \\
\hline
Карточка сделки & Просмотр параметров, истории статусов, платежей и комментариев & Трейдер, позиционер, аналитик \\
\hline
Создание сделки & Ввод валютной пары, суммы, курса, контрагента и типа сделки & Трейдер \\
\hline
Согласование & Одобрение, возврат, отклонение, ввод комментария & Позиционер \\
\hline
Отчеты & Формирование и выгрузка отчетов по заданным фильтрам & Аналитик, руководитель \\
\hline
Администрирование & Пользователи, роли, настройки, справочники & Администратор \\
\hline
\end{tabularx}
\end{table}

\newpage
\begin{center}
\textbf{Приложение В. Логическая модель данных}
\end{center}

Логическая модель данных Системы включает следующие основные сущности:
\begin{itemize}[itemsep=1mm]
  \item \texttt{Deal} -- сделка;
  \item \texttt{Payment} -- расчетный платеж;
  \item \texttt{Currency} -- валюта;
  \item \texttt{Counterparty} -- контрагент;
  \item \texttt{NostroAccount} -- ностро-счет;
  \item \texttt{BusinessCalendar} -- производственный календарь;
  \item \texttt{User} -- пользователь;
  \item \texttt{Role} -- роль;
  \item \texttt{AuditEvent} -- событие аудита.
\end{itemize}

\begin{table}[H]
\centering
\small
\caption{Основные связи логической модели}
\begin{tabularx}{\linewidth}{|L{40mm}|L{40mm}|X|}
\hline
\textbf{Сущность 1} & \textbf{Сущность 2} & \textbf{Связь} \\
\hline
\texttt{Deal} & \texttt{Payment} & Одна сделка содержит два и более расчетных платежа. \\
\hline
\texttt{Deal} & \texttt{Counterparty} & Каждая сделка связана с одним контрагентом. \\
\hline
\texttt{Payment} & \texttt{Currency} & Каждый платеж выполняется в одной валюте. \\
\hline
\texttt{Payment} & \texttt{NostroAccount} & Платеж может быть связан с расчетным счетом ностро. \\
\hline
\texttt{User} & \texttt{Role} & Пользователь может иметь одну или несколько ролей. \\
\hline
\texttt{AuditEvent} & \texttt{User} & Событие аудита фиксирует пользователя, выполнившего действие. \\
\hline
\end{tabularx}
\end{table}

\newpage
\begin{center}
\textbf{Приложение Г. Модель угроз и нарушителя информационной безопасности}
\end{center}

Модель угроз и нарушителя для Системы подлежит уточнению при разработке частной документации по информационной безопасности. На уровне технического проекта выделяются следующие укрупненные категории нарушителей:
\begin{itemize}[itemsep=1mm]
  \item внешний нарушитель, не имеющий легитимного доступа к Системе;
  \item внутренний пользователь с недостаточными полномочиями;
  \item пользователь с легитимными полномочиями, выполняющий ошибочные или вредоносные действия;
  \item технический нарушитель, воздействующий на каналы интеграции или инфраструктуру.
\end{itemize}

Перечень угроз сформирован на основании реестра рисков проекта. Для каждой угрозы указаны связанная уязвимость, актив и проектная мера снижения риска.

{\centering
\footnotesize
\setlength{\tabcolsep}{2pt}
\begin{longtable}{|C{7mm}|L{36mm}|L{40mm}|L{30mm}|L{34mm}|}
\caption{Угрозы информационной безопасности по реестру рисков}\label{tab:threat-register}\\
\hline
\textbf{№} & \textbf{Угроза} & \textbf{Уязвимость} & \textbf{Актив} & \textbf{Мера защиты} \\
\hline
\endfirsthead
\hline
\textbf{№} & \textbf{Угроза} & \textbf{Уязвимость} & \textbf{Актив} & \textbf{Мера защиты} \\
\hline
\endhead
\hline
\multicolumn{5}{|r|}{\textit{Продолжение на след. странице}} \\
\hline
\endfoot
\hline
\endlastfoot

1 & Подмена бизнес-логики (модификация кода сервиса) & Использование динамической загрузки модулей без проверки целостности & Deal Service, Payment Service, Audit Service & Контроль целостности модулей, запрет неподписанных расширений, проверка артефактов сборки. \\
\hline
2 & Несанкционированный вызов методов создания и редактирования сделок & Отсутствие валидации входной JSON-схемы & Payment Controller, Deal Controller & Валидация схем запросов, авторизация методов API, проверка допустимых переходов статусов. \\
\hline
3 & SQL-инъекция через небезопасную конкатенацию запросов & Использование сырых SQL-запросов вместо ORM с параметризацией & Deal Repository, Payment Repository, Audit Repository & Параметризованные запросы, ORM/Query Builder, статический анализ и тестирование на SQLi. \\
\hline
4 & Перехват данных кэша при передаче & Отсутствие пароля на Redis & Кэш Redis & Аутентификация Redis, защищенный сетевой контур, TLS для внутреннего обмена. \\
\hline
5 & Модификация платежных данных в обход бизнес-логики (прямое изменение БД) & Отсутствие шифрования чувствительных полей в БД (номера счетов) & Payment Repository & Ограничение прямого доступа к БД, аудит изменений, шифрование чувствительных атрибутов. \\
\hline
6 & Несанкционированное чтение данных (утечка сделок, позиций, аудита) & Слабые учетные данные суперпользователя PostgreSQL & База данных PostgreSQL & Политика сложных паролей, запрет общего суперпользователя, ротация секретов, контроль доступа. \\
\hline
7 & Удаление старых записей аудита без ведома службы ИБ & Нет механизма неизменяемых логов (append-only) & Audit Service & Неизменяемое хранение аудита, разграничение прав удаления, внешняя архивация журналов. \\
\hline
8 & Подмена образа контейнера при использовании Docker & Недостаточная валидация входных параметров (SQLi, Command Injection) & Backend-приложение API & Подписание и проверка образов, сканирование контейнеров, контроль командных параметров. \\
\hline
9 & Физическая кража сервера или диска & Отсутствие резервирования (RAID, второй БП, горячая замена дисков) & Техническая инфраструктура & Физическая защита ЦОД, шифрование дисков, резервирование и инвентарный контроль носителей. \\
\hline
10 & Подмена JavaScript-кода (загрузка вредоносного фронтенда) & Отсутствие Content Security Policy (CSP) & Веб-интерфейс Frontend & Настройка CSP, контроль целостности статических ресурсов, защищенная поставка фронтенда. \\
\hline
11 & Несанкционированное копирование базы данных сделок & Недостаточная защита дампов БД (слабые пароли архивов) & Информация (данные) & Шифрование дампов, ограничение доступа к архивам, журналирование операций резервного копирования. \\
\hline
12 & Перехват отчетных данных (конфиденциальность сделок) & Передача отчетов по незащищенному каналу (HTTP вместо HTTPS/TLS) & Система отчетности & Обязательное применение HTTPS/TLS, запрет открытых каналов, контроль сертификатов. \\
\hline
13 & Подмена справочных данных (валют, контрагентов, календаря) в канале связи & Нет проверки контрольных сумм или электронной подписи полученных справочников & Поставщик НСИ & Проверка подписи и контрольных сумм, версионирование справочников, журналирование загрузок. \\
\hline
14 & Перехват или подмена данных при передаче (REST/сообщения) & Отсутствие взаимной аутентификации (mTLS) & Система ведения позиций & mTLS между системами, идентификаторы корреляции, повторная проверка источника сообщений. \\
\hline
15 & Сброс пароля BIOS/ОС при физическом доступе & Отсутствие разделения ролей администратора и аудитора & Работник: администратор & Разделение административных и аудиторских полномочий, контроль физического доступа, учет действий. \\
\hline
16 & Искажение или пропуск ошибок в сделках & Слабые парольные политики & Работник: позиционер & Усиление парольной политики, обучение позиционеров, контроль решений и выборочный аудит сделок. \\
\hline
17 & Несанкционированный доступ к критическим операциям подтверждения & Недостаточный аудит действий подтверждения & Работник: позиционер & Детальный аудит подтверждений, фиксация комментариев, контроль аномальных действий. \\
\hline
18 & Повышение привилегий (подтверждение без проверки) & Отсутствие разделения обязанностей: один пользователь может создать и подтвердить сделку & Работник: позиционер & Разделение ролей создания и подтверждения, запрет самоподтверждения, маршрутизация согласования. \\
\hline
19 & Использование легитимных механизмов для несанкционированных действий (злоупотребление полномочиями) & Неформализованные бизнес-правила валидации сделок & Работник: трейдер & Формализация бизнес-правил, автоматические проверки, журналирование действий трейдера. \\
\hline
20 & Социальная инженерия & Отсутствие обязательной двухфакторной аутентификации & Работник: трейдер & Введение двухфакторной аутентификации, обучение пользователей, контроль подозрительных входов. \\
\hline
\end{longtable}
\par}
\normalsize

\newpage
\begin{center}
\textbf{Список изменений}
\end{center}

\begin{table}[H]
\centering
\small
\begin{tabularx}{\linewidth}{|L{25mm}|L{25mm}|X|L{30mm}|}
\hline
\textbf{Дата} & \textbf{Версия} & \textbf{Описание изменений} & \textbf{Автор} \\
\hline
\_\_.\_\_.\DocYear & \DocVersion & Первичная редакция пояснительной записки к техническому проекту. & \TBD \\
\hline
\end{tabularx}
\end{table}

\end{document}
